\begin{abstract}		
\noindent
Original Pennsylvania Motorcycle Operator Manual (PUB.147) PDF
%	 from State Library of Pennsylvania 
\url{https://digitalcollections.statelibrary.pa.gov/Documents/Detail/motorcycle-operator-manual-pub.-147/465384}
	
	\begin{description}
		\item[Title] Motorcycle Operator Manual, Pub. 147
		\item[Date] March 2017
%		\item[Creator] Department of Transportation
%		\item[Subject] Automobile drivers' tests - Handbooks, manuals, etc.
		\item[Rights Details] No Copyright - United States
	\end{description}
	
	\begin{rus}

	Документ переведен с английского на русский максимально формально и близко к оригиналу. Перевод идет бок о бок с оригиналом. Читаем по-русски, тут же перечитываем оригинал, т.к. в любом случае надо знать термины по-английски.
	
	Документ очень хороший, в нем кратко собрано все, что нужно знать для езды на мотоцикле. Т.к. подача сильно сжатая "--- вникаем в каждое слово.
	
	
	Кое-где сноски от переводчика, которых не было в оригинале.
	Предполагается, что у соискателя категории М есть категория С.
	
	Если не терпится сдать теорию, под честное слово вернуться и прочитать все "--- в самом конце ответы с объяснениями. Достаточно прочитать их раз 10 чтобы тупо сдать экзамен. Но этого не достаточно чтобы быть хоть сколько-то готовым к BRC и тем более к езде по дорогам.
	
	После сдачи теории в DMV и сдачи теории в рамках курса BRC, чтобы не засорять голову, надо прочитать книгу «Искусство вождения мотоцикла» Дэвида Хафа (Proficient Motorcycling: The Ultimate Guide to Riding Well. By David L. Hough). Искать старый перевод, с оранжевым KTM на обложке, потому, что новый в синей обложке переведен так ужасно, что местами смысл фраз меняется на противоположный (Мотоцикл. Полное руководство по мастерству вождения. Эксмо - не покупать!).
	
	Переведено с помощью DeepL.com, Proton Lumo, Google Gemini Pro. 
	
%	Сразу хочется сказать про шлем прочую защиту. Особенно в районе Филадельфии легко себя обнаружить за пределы штата. Так что, как говорится, если ты в Риме, то веди себя как римлянин. Как бы Пенсильвания ни была расслаблена в плане правил, про баловство со шлемами надо забыть: никаких скорлупок, фриц-касок, модулей или 3/4 "--- только полноценный литой фулфейс. И устойчивая к истиранию спец. одежда.
	
%	Езда в междурядье: может виной всему соц. сети и видосики без геометок "--- в PA это под запретом. Точка.



%	\bigskip
%	Возможно нужно пояснение к переводу. терминология (rider vs driver, cushion) и казусы всякие.
%	\bigskip
%	
%	
%	\bigskip
	\noindent
	На момент перевода процедура открытия категории М такая: 
	\begin{enumerate}[noitemsep]
		\item Учим теорию
		\item Тренируемся отвечать на вопросы (тут и много онлайн вариантов)
		\item Получаем в DMV лернинг пермит (\$12, проверка зрения, тест теории на компьютере)
		\item Записываемся и идем в школу BRC (в Пенсильвании бесплатно!)
		\item Школа ставит печать, можно ездить на мото полноценно
		\item Ждем письма из DMV, получив идем за новым пластиком
	\end{enumerate}
	
	\noindent
	Исходники в формате \LaTeX{} \\ 
	\link{https://github.com/1001bug/PA-Motorcycle-Operator-Manual}{GitHub Repository}
	
	
	\noindent
	Онлайн тестирование \\ 
	\link{https://1001bug.github.io/PA-Motorcycle-Operator-Manual/online-test/}{Knowledge test EN}
	\noindent\\
	\link{https://1001bug.github.io/PA-Motorcycle-Operator-Manual/online-test/?ru}{Онлайн тест на руссом}
	
	\end{rus}
\end{abstract}




% термины с проблемой перевода
% share the road - делить дорогу, взаимная вежливость на дороге» / «взаимоуважение участников движения», уступайте друг другу дорогу
% rider vs driver - лакончиное rider превращается в длинного мотоциклиста
%